\chapter{Conclusiones}
\label{ch.ej5:conclusiones}

El desarrollo del Ejercicio Práctico N° 5 permitió comprender experimentalmente la configuración y operación de la
generación de señales PWM por hardware utilizando los temporizadores de la arquitectura STM32F103C8T6.

\section{Dificultades Encontradas}
Durante el proceso de diseño e implementación en BareMetal con CMSIS, se presentaron las siguientes dificultades:

\begin{enumerate}
    \item \textbf{Habilitación de Salida Principal en Temporizadores Avanzados (\texttt{TIM1}):} A diferencia de los
    temporizadores de propósito general (\texttt{TIM2}--\texttt{TIM5}), el temporizador avanzado \texttt{TIM1} requiere
    activar obligatoriamente el bit \lstinline{MOE} (\textit{Main Output Enable}) en el registro \lstinline{BDTR}. En las
    primeras pruebas, aunque el canal y el contador estaban activos, no había señal física en el pin \texttt{PA8}
    hasta configurar dicho registro.
    \item \textbf{Configuración del Módulo GPIO en Función Alternativa:} Resultó imprescindible configurar el pin
    \texttt{PA8} como salida de \textit{Función Alternativa Push-Pull} (\lstinline{CNF8 = 10b}). Una configuración
    como salida push-pull genérica desconecta internamente la salida del periférico \texttt{TIM1} de la línea física.
    \item \textbf{Sincronización mediante Registros Sombra:} Se comprendió la necesidad de habilitar los bits de
    precarga (\lstinline{OC1PE} y \lstinline{ARPE}) para evitar anomalías temporales (\textit{glitches}) al actualizar
    el registro \lstinline{CCR1} durante la ejecución del bucle principal.
\end{enumerate}

\section{Conceptos Aprendidos}
Entre los conceptos técnicos consolidados a lo largo de la práctica se destacan:

\begin{enumerate}
    \item \textbf{Generación Eficiente de PWM por Hardware:} La delegación de la conmutación de señales periódicas al
    periférico del temporizador libera totalmente al procesador ARM Cortex-M3, permitiéndole ejecutar otras tareas en
    paralelo de forma concurrente.
    \item \textbf{Ajuste Fino de Frecuencia y Resolución:} Se dominó el cálculo que vincula la frecuencia de reloj
    ($f_{\mathrm{CK\_PSC}}$), el preescalador (\texttt{PSC}) y el registro de autorrecarga (\texttt{ARR}) para fijar una
    frecuencia de PWM deseada ($\qty{1}{\kHz}$) preservando una resolución de 1000 pasos discretos para el ciclo de trabajo.
    \item \textbf{Arquitectura Concurrente No Bloqueante:} Se logró integrar la base de tiempo determinista obtenida con
    \textit{SysTick} para realizar el parpadeo testigo del LED en \texttt{PC13} al mismo tiempo que se efectuaba la rampa
    de variación de brillo en el LED de potencia conectado a \texttt{PA8}.
\end{enumerate}
